% Unit shorthands, \input by the thesis and by paper-notes. Requires siunitx.
%
% \keV{511}, not \qty{511}{\kilo\electronvolt}. Add a line here the first time
% you need a new unit. For a range, use siunitx's \qtyrange directly --
% \qtyrange{4}{12}{\nano\second}.

\NewDocumentCommand\keV{m}{\qty{#1}{\kilo\electronvolt}}
\NewDocumentCommand\MeV{m}{\qty{#1}{\mega\electronvolt}}
\NewDocumentCommand\ps{m}{\qty{#1}{\pico\second}}
\NewDocumentCommand\ns{m}{\qty{#1}{\nano\second}}
\NewDocumentCommand\mm{m}{\qty{#1}{\milli\metre}}
\NewDocumentCommand\cm{m}{\qty{#1}{\centi\metre}}
\NewDocumentCommand\um{m}{\qty{#1}{\micro\metre}}
\NewDocumentCommand\mins{m}{\qty{#1}{\minute}} % not \minute -- that is siunitx's own unit macro
\NewDocumentCommand\degs{m}{\qty{#1}{\degree}} % not \deg -- that is the LaTeX math operator
\NewDocumentCommand\MBq{m}{\qty{#1}{\mega\becquerel}}

\DeclareSIUnit\curie{Ci} % non-SI, so siunitx does not predefine it
